\section{Proofs for the stable function-domain construction}
\label{annex:p10-stable-domains-sequential-models}

\begin{annexitem}{annex:p10.sd.foundations}{thm:p10-sd-foundations}{Finite compact ideals and the required lattice operations}
\begin{proof}
If \(x\sqsubseteq a\) with \(a\) compact, the compact approximants
below \(x\) form a directed set inside the finite \(\downarrow a\).
A directed subset of a finite poset has a greatest member, by finitely
many upper-bound choices within it. Its supremum \(x\) is therefore
one of these compact approximants. This proves that every point in
\(\downarrow a\) is compact. A finite bounded join of compact points
is compact: if below a directed supremum, compactness supplies an index
above each constituent and directedness gives one index above all, hence
above their join.

For any nonempty subset \(X\), its lower bounds form a subset bounded
above by any chosen member of \(X\), and include bottom. Their bounded
join is the greatest lower bound of \(X\). A subset of \(\downarrow u\)
has its bounded join there, including the empty join; its infimum there
is the preceding nonempty infimum, or \(u\) for the empty subset.
Thus the principal ideal is a complete lattice. The stated distributive
law restricts to it.
To prove joint continuity of binary meet, let \((x_i,y_i)\) be directed,
and let compact \(c\) lie below \((\bigsqcup x_i)\sqcap(\bigsqcup y_i)\).
Compactness gives indices with \(c\) below each coordinate; a common
upper index gives \(c\sqsubseteq x_j\sqcap y_j\). Algebraicity of
the meet element proves it below \(\bigsqcup_i(x_i\sqcap y_i)\).
The other inequality is monotonicity. Hence meet is jointly continuous.

Finite products have coordinatewise directed joins, bounded joins and
nonempty meets. Compact tuples are exactly tuples of compact coordinates:
the forward direction follows by varying one coordinate through its
compact approximation while holding the others fixed; the reverse uses
finitely many compactness witnesses and one directed upper index.
These tuples approximate every product point. Their principal ideals are
finite products of finite principal ideals, and distributivity is
coordinatewise. The empty product is the one-point dI-domain.
\end{proof}
\end{annexitem}

\begin{annexitem}{annex:p10.sd.moduli}{thm:p10-sd-moduli}{Least supports and the source's stronger function order}
\begin{proof}
Let \(f\) be stable and \(b\in K(E)\) lie below \(f(x)\).
Continuity and algebraicity give compact \(a\sqsubseteq x\) with
\(b\sqsubseteq f(a)\). Among the finite \(\downarrow a\), take
all witnesses \(d\) with \(b\sqsubseteq f(d)\). They are nonempty
and have common bound \(a\). Their finite meet \(m\) is a witness
by repeated stability. For any other witness \(w\sqsubseteq x\),
\(a,w\) have common bound \(x\), so \(a\sqcap w\) is a witness
in \(\downarrow a\). Consequently \(m\sqsubseteq a\sqcap w\sqsubseteq w\).
Thus \(m\) is the least witness below \(x\), not merely a minimal
one in the finite interval; it is compact by the preceding theorem.

Conversely assume all moduli exist. If \(x,y\sqsubseteq z\) and
compact \(b\sqsubseteq f(x)\sqcap f(y)\), then \(M(f,z,b)\) lies
below both witnesses \(x,y\). It lies below \(x\sqcap y\), so
\(b\sqsubseteq f(x\sqcap y)\). Algebraicity proves the meet inequality
for all output information; the reverse inequality follows from
monotonicity. This establishes stability. For any nonempty bounded
family \(X\) with bound \(z\), the same argument for compact
\(b\sqsubseteq\bigsqcap f(X)\) puts its modulus below every member
and hence below \(\bigsqcap X\). Algebraicity and monotonicity give
preservation of its infimum.

If \(f\sqsubseteq_s g\), setting \(x=y\) gives pointwise
\(f\sqsubseteq g\). For compact \(b\sqsubseteq f(x)\), put
\(m_f=M(f,x,b)\), \(m_g=M(g,x,b)\). The pointwise inequality gives
\(m_g\sqsubseteq m_f\). Since \(m_g\sqsubseteq x\), the stable-order
equation gives \(f(m_g)=f(x)\sqcap g(m_g)\), which contains \(b\).
Hence \(m_f\sqsubseteq m_g\), proving equality.
Conversely assume pointwise order and modulus agreement. For
\(x\sqsubseteq y\) and compact \(b\sqsubseteq f(y)\sqcap g(x)\),
the modulus for \(g\) below \(y\) is below its witness \(x\).
The equal \(f\)-modulus therefore witnesses \(b\sqsubseteq f(x)\).
Algebraicity gives \(f(y)\sqcap g(x)\sqsubseteq f(x)\); the other
inequality follows from monotonicity and pointwise order. This is exactly
the stable-order equation.
\end{proof}
\end{annexitem}

\begin{annexitem}{annex:p10.sd.ordered-functions}{thm:p10-sd-ordered-functions}{Directed completeness and stable bounded joins without a trace assumption}
\begin{proof}
Reflexivity uses monotonicity of a stable map. The order equation at equal
inputs gives pointwise order, so antisymmetry follows. If
\(f\sqsubseteq_s g\sqsubseteq_s h\) and \(x\sqsubseteq y\), then
\(f(y)\sqcap h(x)=f(y)\sqcap g(y)\sqcap h(x)=f(y)\sqcap g(x)=f(x)\),
proving transitivity. Bottom satisfies the equation with every map.

For a directed stable-ordered family \((f_i)\), let \(f\) be its
pointwise supremum, which is continuous by the ordinary function-dcpo
theorem. Joint continuity of meet in \(E\) gives, for compatible inputs,
\(f(x)\sqcap f(y)=\bigsqcup_i(f_i(x)\sqcap f_i(y))=f(x\sqcap y)\).
Thus it is stable. To show \(f_i\sqsubseteq_s f\), for any index
\(j\) choose \(k\) above \(i,j\). At \(x\sqsubseteq y\),
\(f_i(y)\sqcap f_j(x)\sqsubseteq f_i(y)\sqcap f_k(x)=f_i(x)\).
The diagonal index supplies equality in the supremum, and continuity
of meet gives the required equation with \(f(x)\). If every \(f_i\)
is below \(g\) stably, take their directed supremum in
\(f_i(y)\sqcap g(x)=f_i(x)\) to prove \(f\sqsubseteq_s g\).
This proves the actual least upper bound, not just a continuous pointwise map.

A useful interval fact is: if \(u,v\sqsubseteq_s h\), then
\(u\sqsubseteq v\) pointwise implies \(u\sqsubseteq_s v\).
Indeed \(u(y)\sqcap v(x)=u(y)\sqcap v(y)\sqcap h(x)=u(y)\sqcap h(x)=u(x)\).
Within this interval stable order is therefore pointwise order.
For \(u,v\sqsubseteq_s h\), define \(w(x)=u(x)\sqcup v(x)\).
The join exists below \(h(x)\). Finite bounded joins commute with
directed coordinate suprema, by their upper-bound definitions; applying
this to a directed input family proves continuity of \(w\).
Bounded distributivity gives
\(w(y)\sqcap h(x)=w(x)\), so \(w\sqsubseteq_s h\).
For compatible \(x,y\) bounded by \(z\), expand
\(w(x)\sqcap w(y)\) into the four cross terms. For any two members
\(u_i,u_j\in\{u,v\}\),
\[
 u_i(x)\sqcap u_j(y)
 =u_i(z)\sqcap u_j(z)\sqcap h(x)\sqcap h(y)
 \sqsubseteq u_i(z)\sqcap h(x\sqcap y)=u_i(x\sqcap y).
\]
The diagonal terms give the reverse inclusion, so \(w\) is stable.
The interval fact puts both \(u,v\) below \(w\) stably. The same
equation holds with any common stable upper bound instead of \(h\),
so \(w\) is their stable least upper bound. Finite joins extend this
argument. The directed family of finite joins of any bounded family
then gives its pointwise stable supremum, including bottom for the empty family.

Pointwise \(m(x)=u(x)\sqcap v(x)\) is continuous and stable and
satisfies \(m\sqsubseteq_s h\). The interval fact makes it the
stable greatest lower bound of \(u,v\). This identification requires
their common stable upper bound; it is not asserted for incompatible maps.
\end{proof}
\end{annexitem}

\begin{annexitem}{annex:p10.sd.function-domain}{thm:p10-sd-function-domain}{Compact finite cuts, algebraicity, property I and distributivity}
\begin{proof}
Meet continuity proves continuity of \(f_{a,b}\). For compatible
\(x,y\), the cut inputs have common bound \(a\); stability of \(f\)
and associativity of meet show preservation of their meet, so the cut
is stable. At \(x\sqsubseteq y\), \(x,y\sqcap a\) have common bound
\(y\), and
\(f_{a,b}(y)\sqcap f(x)=f(y\sqcap a)\sqcap b\sqcap f(x)
=f(x\sqcap a)\sqcap b=f_{a,b}(x)\).
Thus the cut lies below \(f\) in the stable order.

To prove compactness of a cut \(c=f_{a,b}\), suppose
\(c\sqsubseteq_s F=\bigsqcup_i f_i\) for a directed stable family.
The finite set \(\downarrow a\) consists of compact points.
For each \(y\) in it, \(c(y)\sqsubseteq b\) is compact. Its inequality
below the pointwise supremum gives an index with \(c(y)\sqsubseteq f_i(y)\).
One common upper index \(i_0\) works for the whole finite set.
For arbitrary \(x\), the cut identity \(c(x)=c(x\sqcap a)\)
and monotonicity give \(c(x)\sqsubseteq f_{i_0}(x)\).
Both \(c,f_{i_0}\) are below \(F\) stably; the interval fact now
upgrades pointwise order to \(c\sqsubseteq_s f_{i_0}\).
This is genuine stable-order compactness. Finite bounded joins of these
compact maps are compact by finitely many directed witnesses.

The finite joins of cuts of \(f\) are a directed stable-bounded family.
Their supremum is pointwise. For each compact \(b\sqsubseteq f(x)\),
algebraicity and continuity supply compact \(a\sqsubseteq x\) with
\(b\sqsubseteq f(a)\). Then \(f_{a,b}(x)=b\), so the supremum
contains every compact approximant to \(f(x)\). It equals \(f\).
Thus the function dcpo is algebraic. If \(f\) itself is compact, it
is below one finite cut join in its own family and hence equals that
join, since every member is below it stably.

Let compact \(h=\bigsqcup_{j=1}^n h_{a_j,b_j}\), using its own cuts,
and set \(Y=\bigcup_j\downarrow a_j\), a finite set. For
\(k\sqsubseteq_s h\), bounded distributivity and its order equation give
\[
 k(x)=\bigsqcup_j\bigl(k(x\sqcap a_j)\sqcap h_{a_j,b_j}(x)\bigr).
\]
Indeed expand \(k(x)\sqcap h(x)\), and note that each cut value lies
below \(h(x\sqcap a_j)\); then use
\(k(x\sqcap a_j)=k(x)\sqcap h(x\sqcap a_j)\).
Hence the values on \(Y\) determine \(k\). For \(y\in Y\), each
cut value below \(b_j\) is compact, so their finite join \(h(y)\)
is compact. Its principal ideal is finite by property I in \(E\).
There are finitely many possible tables \((k(y))_{y\in Y}\), proving
that \(\downarrow_s h\) is finite.

Bounded completeness was already proved. Its nonempty infima exist by
taking joins of lower bounds. In an interval below \(p\), binary stable
meets and bounded joins are pointwise and are distributive by the law
in \(E\). For arbitrary \(x\) and compatible \(y,z\sqsubseteq_s p\),
replace \(x\) by \(x\sqcap_s p\); its meets with \(y,z,y\sqcup_s z\)
are unchanged. All arguments then lie in that interval, proving the
global bounded distributive law. This completes the dI construction.
For countable compact bases, a cut with fixed \(a,b\) is determined
by a finite table from \(\downarrow a\) to \(\downarrow b\).
Only finitely many such tables are possible. Countably many \(a,b\)
and finite joins therefore enumerate a countable set of compact maps.
No effective decision procedure for those tables is asserted.
\end{proof}
\end{annexitem}

\begin{annexitem}{annex:p10.sd.cartesian-closed}{thm:p10-sd-cartesian-closed}{Stable evaluation and the actual exponential bijection}
\begin{proof}
Identity is stable and composition preserves continuity and compatible
meets; a compatible input pair has compatible images by monotonicity.
Product maps and pairing are stable coordinatewise; the one-point object
is terminal. Evaluation is continuous because directed stable function
suprema are pointwise, and ordinary continuous evaluation preserves the
directed diagonal supremum of a pair of function/input families.
For compatible \((f,x),(g,y)\) with common bound \((h,z)\), their
function meet is pointwise. The stable-order equations give
\[
 f(x)\sqcap g(y)=f(z)\sqcap g(z)\sqcap h(x)\sqcap h(y)
 =f(z)\sqcap g(z)\sqcap h(x\sqcap y)
 =f(x\sqcap y)\sqcap g(x\sqcap y).
\]
This is evaluation at the product meet, proving joint stability.

For stable \(k:C\times D\to E\), define \(\Lambda k(c)(d)=k(c,d)\).
The slice in \(d\) is stable. If \(c\sqsubseteq c'\) and
\(d\sqsubseteq d'\), the two product inputs \((c,d'),(c',d)\)
have common bound \((c',d')\), so
\(k(c,d)=k(c,d')\sqcap k(c',d)\).
This is exactly \(\Lambda k(c)\sqsubseteq_s\Lambda k(c')\).
Continuity in \(c\) follows pointwise from that of \(k\) and the
function-supremum theorem. For compatible \(c,c'\), their slices
have a common stable upper slice. At each \(d\), stability of \(k\)
gives \(k(c\sqcap c',d)=k(c,d)\sqcap k(c',d)\); their stable
meet is pointwise, proving stability of the curried map. Conversely,
uncurrying is stable evaluation after the product map; the two procedures
are inverse by their values and natural by substitution/composition.

For the hom-order assertion, suppose \(f\sqsubseteq_s f':D\to E\)
and \(g\sqsubseteq_s g':E\to F\). At \(x\sqsubseteq y\),
\(f(y),f'(x)\) have common bound \(f'(y)\). Using both stable-order
equations and stability of \(g'\),
\[
 g(f(y))\sqcap g'(f'(x))
 =g(f'(y))\sqcap g'(f(y))\sqcap g'(f'(x))
 =g(f'(y))\sqcap g'(f(x))=g(f(x)).
\]
Thus composition is monotone in stable hom orders. Directed hom suprema
are pointwise; continuity of the outer maps and directed diagonal
cofinality give \((\bigsqcup_i g_i)(\bigsqcup_i f_i)(x)
=\bigsqcup_i g_i(f_i(x))\). Independent directed families use their
directed product of indices. This proves local continuity as claimed.
\end{proof}
\end{annexitem}

\begin{annexitem}{annex:p10.sd.stable-fix}{thm:p10-sd-stable-fix}{The common-upper-bound chain makes fixed points stable}
\begin{proof}
Stable endomaps are continuous on the pointed dcpo, so their ordinary
least fixed points exist. Stable order implies pointwise order, hence
iteration makes fix monotone. The maps \(f\mapsto f^n(\bot)\) are
continuous by induction using jointly continuous stable evaluation.
Their increasing pointwise supremum is fix. Interchanging the directed
suprema over the map family and over \(n\) proves Scott continuity.

Suppose \(f,g\sqsubseteq_s h\). Their stable meet \(k\) is pointwise.
Put \(a_n=f^n\bot,b_n=g^n\bot,c_n=h^n\bot\); monotonicity of
iteration gives \(a_n,b_n\sqsubseteq c_n\). We prove
\(k^n\bot=a_n\sqcap b_n\). The base is bottom. At the induction
step, stable-order equations below \(c_n\) and stability of \(h\) give
\[
 \begin{split}
 k(a_n\sqcap b_n)
 &=f(c_n)\sqcap g(c_n)\sqcap h(a_n\sqcap b_n)\\
 &=f(c_n)\sqcap g(c_n)\sqcap h(a_n)\sqcap h(b_n)\\
 &=f(a_n)\sqcap g(b_n)=a_{n+1}\sqcap b_{n+1}.
 \end{split}
\]
The chains have the stated common bounds at every step; no unbounded
meet-preservation premise is smuggled in. Joint continuity of meet in
\(D\) takes their directed limits to
\(\operatorname{fix}(k)=\operatorname{fix}(f)\sqcap\operatorname{fix}(g)\).
Together with continuity this proves that fix itself is a stable map.
\end{proof}
\end{annexitem}

\begin{annexitem}{annex:p10.sd.pcf-interpretation}{thm:p10-sd-pcf-interpretation}{Every PCF constructor, including parameterized stable recursion}
\begin{proof}
Flat ground domains satisfy all dI conditions; products and the proved
stable function constructor therefore build \(S_t\) by type induction.
Interpret terms by typing induction. Projections and constants are stable.
Each strict ground unary operation is continuous and preserves compatible
meets in the flat input: compatible nonbottom inputs coincide, and a
bottom input has bottom output. For the conditional on the product of
guard and two branches, compatible nonbottom guards coincide and choose
the same branch, whose meet is the output meet. If either guard is
bottom, the product-meet guard and one output are bottom, giving equality.
Its continuity is proved by guard stabilization on a cofinal tail exactly
as for the continuous model; no branch is forced unnecessarily.
Application uses stable evaluation, and abstraction stable currying.
The recursion body is stable on the environment product with its bound
variable by the induction hypothesis. Currying gives a stable map from
environments into stable endomaps; composing with the proved stable fix
gives stable dependence of the recursive meaning. This closes the
source's Lemma~4.4.5 dependency at arbitrary simple result types.

Semantic substitution follows by syntax induction: variables select the
supplied meanings; constants and each operation/application combine the
induction hypotheses. Fresh binder names make updated environments
commute with substitution. At lambda, equality holds at every argument;
at mu, the body equality holds for every bound-variable value, making
the two defining endomaps and their least fixed points equal.
This also proves bound-renaming invariance. Beta uses substitution and
evaluation; mu unfolding uses the least-fixed-point equation; ground and
conditional contractions use the same declared values, including \(\operatorname{pred}(0)=0\).
Compositionality propagates equality through each evaluation context.
Hence the interpretation is both well defined and operationally sound.
\end{proof}
\end{annexitem}

\begin{annexitem}{annex:p10.sd.pcf-adequacy}{thm:p10-sd-pcf-adequacy}{The logical relation is rebuilt on the stable carriers}
\begin{proof}
Define \(a\mathrel{\mathcal R^s_t}M\) on \(S_t\) and closed source
PCF terms: bottom is related to every ground term, nonbottom ground
points require the corresponding evaluation, and a stable function
is related at \(s\to t\) iff it sends every related argument
element/closed term pair to a related result with syntax \(M N\).
The fixed-term fibers contain bottom and are directed-closed by type
induction. Ground flat compactness extracts a nonbottom supremum from
one member. At arrows a directed stable supremum is pointwise, so
evaluation at each related argument reduces closure to the result type.
Small-step conversion preserves and reflects this relation: the ground
case uses the proved source evaluation correspondence; at arrows operator
context propagation and result-type induction handle \(M N\to M' N\).
Thus the relation is justified on these carriers, not identified without
proof with a relation whose functions live in the full continuous model.

The fundamental lemma now follows by typing induction for related
environments/substitutions. Variables and constants are immediate.
Application is the arrow clause. At lambda extend the related assignments
by an arbitrary related argument pair, use the body hypothesis, and
convert backwards through beta. A bottom numerical input yields bottom
under every strict unary constructor; a nonbottom input's relatedness
provides its numeral evaluation and the arithmetic rule. A bottom guard
yields bottom at any result type; a nonbottom guard evaluates to its
boolean, and the selected branch hypothesis plus context conversion
gives relatedness of the conditional.
For mu, with closed recursive syntax \(Q\) and stable continuous body
endomap \(F\), prove \(F^n\bot\mathrel{\mathcal R^s_t}Q\) by
numerical induction. Bottom gives the base. The smaller body's typing
hypothesis under the assignment \(x\mapsto F^n\bot,Q\) gives
relatedness to the unfolded syntax; backwards unfolding conversion gives
the next iterate related to \(Q\). Directed closure gives the least
fixed point related to \(Q\). This proves the entire fundamental lemma
with the actual stable model's recursion and suprema.

In the empty context a program's own denotation is related to it.
A nonbottom ground meaning therefore supplies its exact value evaluation.
Conversely operational soundness gives that meaning from every finite
evaluation. The flat ground carriers are exhausted by bottom and these
distinct value meanings, so no evaluation is equivalent to bottom.
For contextual equality, induction on the typed context uses the stable
interpretation's substitution/composition and binders to preserve equality;
closing at ground and applying both observation directions yields equal
results or divergence. The original operational progress theorem remains
about the unchanged syntax/strategy. This proves stable-model adequacy,
not a full-abstraction converse.
\end{proof}
\end{annexitem}

\begin{annexitem}{annex:p10.sd.counterexample}{thm:p10-sd-counterexample}{The source's stable separator and exact contextual consequence}
\begin{proof}
The two maps \(p,q\) are continuous and stable: on flat compatible
inputs either the nonbottom values coincide or one input is bottom;
the constant map preserves the meet and the strict map gives bottom in
the latter case. Pointwise \(q\sqsubseteq p\). But at
\(\bot\sqsubseteq\mathsf{true}\),
\(q(\mathsf{true})\sqcap p(\bot)=\mathsf{true}\ne q(\bot)\),
so \(q\not\sqsubseteq_s p\), while \(p\not\sqsubseteq_s q\)
already fails pointwise at bottom.
If \(h\) is above \(p\) stably, pointwise order forces \(h(\bot)\)
true and hence all its outputs true; thus \(h=p\).
If \(h\) is above \(q\) stably, both nonbottom outputs must be true.
Its bottom output is bottom or true by monotonicity. True would make
\(h=p\), which is impossible by the just-proved order test. Hence
\(h=q\). Thus \(p,q\) are distinct maximal elements and cannot
have a common stable upper bound.
Any \(k\sqsubseteq_s p\) has values bottom or true. Its order equation
at \(\bot\sqsubseteq y\) is \(k(\bot)=k(y)\sqcap\mathsf{true}=k(y)\),
so it is constant bottom or \(p\). Since \(p\) is not below \(q\),
their only common stable lower bound is bottom. Their pointwise meet is
\(q\), as their pointwise comparison already showed.

The stable function carrier on the finite boolean domain is finite.
The map \(r\), nonbottom only at its two specified maximal points,
is monotone, and every monotone map on a finite dcpo is continuous.
To check stability, take a compatible pair \(u,v\). If both have
nonbottom \(r\)-values, compatibility forces them to be the same one
of \(p,q\), and the meet equation holds. Otherwise the output meet
is bottom. The input meet cannot be \(p\) or \(q\): maximality would
force both inputs equal to that point and make both outputs nonbottom.
Therefore its \(r\)-value is bottom too. This proves actual stable
membership of the semantic separator, not merely its monotonicity.

The displayed closed PCF terms for \(p,q\) have exactly those meanings
in both models, because the conditional is strict in its guard even
with equal branches. In the continuous model, monotonicity of every
semantic \(f\) and \(q\sqsubseteq p\) imply that \(f(q)=\mathsf{true}\)
forces \(f(p)=\mathsf{true}\). Thus the nested terms \(M,N\) both
give bottom on every continuous argument: a bottom outer guard is
bottom, a false guard selects Omega, and a true guard forces the inner
true branch Omega. Their continuous denotations are equal. The already
proved continuous contextual adequacy theorem therefore gives
\(M\simeq_o N\) for the unchanged PCF and all closing program contexts.
In the stable model the separator has \(r(q)=\mathsf{true}\),
\(r(p)=\mathsf{false}\), so the same source terms give true and false,
respectively. This proves failure of stable full abstraction.
If a closed PCF term denoted \(r\), the closing context applying its
hole to that term would make \(M,N\) evaluate to true and false by
stable ground adequacy. That contradicts their contextual equivalence.
Hence the separator is not definable in this language, even though it
is a stable continuous semantic function.

Finally the parallel-or table has value true on both
\((\mathsf{true},\bot)\) and \((\bot,\mathsf{true})\), a compatible
pair bounded by \((\mathsf{true},\mathsf{true})\). Their meet is
\((\bot,\bot)\), where the table gives bottom, not the meet true of
their outputs. Hence it is not stable. Excluding this earlier separator
does not remove the distinct stable separator \(r\).
\end{proof}
\end{annexitem}
